\begin{definition}[5.6.2]
\label{IV.2.s5.6.2-fr}
Se dice que un anillo noetheriano $A$ es universalmente catenario si
satisface las condiciones equivalentes \textup{a), b), c)} de \textup{(5.6.1)}.
\end{definition}

\begin{remarks}[5.6.3]
\label{IV.2.s5.6.3-fr}
\begin{enumerate}
\item[(i)] Si $A$ es universalmente catenario, lo mismo ocurre con
$S^{-1}A$ para toda parte multiplicativa $S$ de $A$, como se sigue
inmediatamente de (5.6.1, \emph{a)}) y del hecho de que todo anillo de fracciones de
un anillo catenario es catenario. Recíprocamente, si, para todo ideal primo
$\mathfrak p$ de $A$, el anillo $A_{\mathfrak p}$ es universalmente
catenario, entonces $A$ es universalmente catenario: esto se sigue de
(5.6.1, \emph{c)}), observando que, poniendo $S=A\setmin\mathfrak p$, $S^{-1}B$ es una
$A_{\mathfrak p}$-álgebra de tipo finito, y
$B_{\mathfrak q}=(S^{-1}B)_{S^{-1}\mathfrak q}$.
\item[(ii)] Se dice que un preesquema localmente noetheriano $X$ es
\emph{universalmente catenario} si, para todo $x\in X$, el anillo
$\mathcal O_x$ es universalmente catenario. Se sigue de (i) que, para que
$A$ sea un anillo universalmente catenario, es necesario y suficiente que lo
sea el esquema $\operatorname{Spec}(A)$.
\item[(iii)] El criterio (5.6.1, \emph{c)}) solo hace intervenir las imágenes de $A$
en $A$-álgebras íntegras de tipo finito, por tanto \emph{anillos cociente
íntegros} de $A$. Se concluye que, si

\oldpage[IV-2]{99}

$\mathfrak p_i$ ($1\leq i\leq n$) son los ideales primos minimales de $A$,
equivale a decir que $A$ es universalmente catenario o que lo es cada uno de
los $A/\mathfrak p_i$. Esto muestra también que equivale a decir que un
preesquema localmente noetheriano $X$ es universalmente catenario, o que lo es
$X_{\mathrm{red}}$.
\item[(iv)] El criterio (5.6.1, \emph{b)}) y la observación (i) muestran que, si $A$
es un anillo universalmente catenario, lo mismo ocurre con toda $A$-\emph{álgebra
esencialmente de tipo finito} (1.3.8).
\end{enumerate}
\end{remarks}

